% The widespread use of Artificial Intelligence (AI) in consequential domains, such as healthcare and parole decision making systems, has drawn intense scrutiny on the fairness of these methods. Fairness by itself is often insufficient as the rationale for a contentious decision often needs to be audited, understood, and defended. 
% We propose that attention mechanism can be used to ensure fair outcomes while simultaneously providing an interpretable accounting for how a decision was made. 
% This can help us to eliminate problematic features and build fair, interpretable, and accurate models. Toward this goal, we design an attention-based model that can be leveraged as an interpretability framework which can identify features responsible for both performance and fairness of the model through attention interventions and attention weight manipulation. Using this interpretability framework, we then design a post-processing bias mitigation strategy and compare it with a suite of baselines. In addition to tabular data, we lastly demonstrate the effectiveness of our model in reducing bias on non-tabular data formats.

% Modified abstract to add "attribution" term and some other minor changes
The widespread use of Artificial Intelligence (AI) in consequential domains, such as healthcare and parole decision-making systems, has drawn intense scrutiny on the fairness of these methods. However, ensuring fairness is often insufficient as the rationale for a contentious decision needs to be audited, understood, and defended. We propose that the attention mechanism can be used to ensure fair outcomes while simultaneously providing feature attributions to account for how a decision was made. Toward this goal, we design an attention-based model that can be leveraged as an attribution framework. It can identify features responsible for both performance and fairness of the model through attention interventions and attention weight manipulation. Using this attribution framework, we then design a post-processing bias mitigation strategy and compare it with a suite of baselines. We demonstrate the versatility of our approach by conducting experiments on two distinct data types, tabular and textual.
%Lastly, in addition to tabular data, we demonstrate the effectiveness of our model in reducing bias on non-tabular data.